\documentclass{article}
\usepackage{amsmath,amssymb,amsthm}
\usepackage{algorithm}
\usepackage{algorithmic}
\usepackage{bm}
\usepackage{hyperref}
\usepackage{enumitem}
\newtheorem{assumption}{Assumption}
\newtheorem{theorem}{Theorem}
\newtheorem{lemma}{Lemma}
\newcommand{\EE}{\mathbb{E}}
\newcommand{\norm}[1]{\left\lVert #1\right\rVert}
\newcommand{\grad}{\nabla}
\newcommand{\prox}{\operatorname{prox}}
\begin{document}

\section*{Algorithm Name}
\textbf{Exponential Moving‑Average Gradient Descent (EMA‑GD)}  

The method keeps an exponential moving average (EMA) of the stochastic gradients and adapts the stepsize according to the magnitude of the bias‑corrected EMA.

-----------------------------------------------------------------------

\section*{Mathematical Formulation}
Let $f:\mathbb{R}^{d}\to\mathbb{R}$ be differentiable.  
Given an initial point $x_{0}\in\mathbb{R}^{d}$, smoothing parameter $\alpha\in(0,1]$, base stepsize $\eta>0$ and scaling constant $\beta\ge 0$, EMA‑GD generates sequences $\{x_{t}\}_{t\ge 0}$, $\{v_{t}\}_{t\ge 0}$ and bias‑corrected averages $\{\hat v_{t}\}_{t\ge 0}$ as follows:

\begin{align}
    &\text{(stochastic gradient)}\qquad g_{t}= \grad f(x_{t};\xi_{t}), \label{eq:stochgrad}\\[4pt]
    &\text{(EMA update)}\qquad 
        v_{t}= (1-\alpha)v_{t-1}+ \alpha g_{t},\qquad v_{-1}=0, \label{eq:ema}\\[4pt]
    &\text{(bias correction)}\qquad 
        \hat v_{t}= \frac{v_{t}}{1-(1-\alpha)^{t+1}}, \label{eq:bias}\\[4pt]
    &\text{(stepsize)}\qquad 
        \eta_{t}= \frac{\eta}{1+\beta \norm{\hat v_{t}}}, \label{eq:stepsize}\\[4pt]
    &\text{(parameter update)}\qquad 
        x_{t+1}= x_{t}-\eta_{t}\,\hat v_{t}. \label{eq:update}
\end{align}

The algorithm stops after a predetermined number $T$ of iterations or when $\norm{\grad f(x_{t})}\le \varepsilon$.

-----------------------------------------------------------------------

\section*{Pseudocode}
\begin{algorithm}[H]
\caption{EMA‑GD}
\begin{algorithmic}[1]
\STATE \textbf{Input:} $x_{0}$, stepsize $\eta>0$, EMA weight $\alpha\in(0,1]$, scaling $\beta\ge0$, max‑iters $T$.
\STATE $v_{-1}\leftarrow 0$.
\FOR{$t=0,\dots,T-1$}
    \STATE Sample a minibatch (or a single sample) $\xi_{t}$ and compute $g_{t}= \grad f(x_{t};\xi_{t})$.
    \STATE $v_{t}\leftarrow (1-\alpha)v_{t-1}+ \alpha g_{t}$ \hfill\COMMENT{EMA of gradients}
    \STATE $\hat v_{t}\leftarrow v_{t}/\bigl(1-(1-\alpha)^{t+1}\bigr)$ \hfill\COMMENT{bias correction}
    \STATE $\eta_{t}\leftarrow \displaystyle\frac{\eta}{1+\beta\norm{\hat v_{t}}}$.
    \STATE $x_{t+1}\leftarrow x_{t}-\eta_{t}\hat v_{t}$.
\ENDFOR
\STATE \textbf{Output:} $x_{T}$ (or the best iterate).
\end{algorithmic}
\end{algorithm}

-----------------------------------------------------------------------

\section*{Dry Run Example}
Consider the one‑dimensional quadratic $f(w)=(w-3)^{2}$, whose gradient is $\grad f(w)=2(w-3)$.  
Take $w_{0}=0$, base stepsize $\eta=0.1$, EMA weight $\alpha=0.5$, and $\beta=0$ (so $\eta_{t}=\eta$).  
We ignore stochasticity and use the true gradient as $g_{t}$.

\begin{center}
\begin{tabular}{c|c|c|c|c}
$t$ & $g_{t}=2(w_{t}-3)$ & $v_{t}$ (EMA) & $\hat v_{t}$ (bias‑corr.) & $w_{t+1}=w_{t}-\eta\hat v_{t}$\\ \hline
0 & $2(0-3)=-6$ &
    $v_{0}=(1-\alpha)0+\alpha(-6)= -3$ &
    $\hat v_{0}= \dfrac{-3}{1-(1-\alpha)^{1}}= \dfrac{-3}{1-0.5}= -6$ &
    $w_{1}=0-0.1(-6)=0.6$\\[6pt]
1 & $2(0.6-3)=-4.8$ &
    $v_{1}=0.5(-3)+0.5(-4.8)= -3.9$ &
    $\hat v_{1}= \dfrac{-3.9}{1-(0.5)^{2}}= \dfrac{-3.9}{1-0.25}= -5.2$ &
    $w_{2}=0.6-0.1(-5.2)=1.12$\\[6pt]
2 & $2(1.12-3)=-3.76$ &
    $v_{2}=0.5(-3.9)+0.5(-3.76)= -3.83$ &
    $\hat v_{2}= \dfrac{-3.83}{1-(0.5)^{3}}= \dfrac{-3.83}{1-0.125}= -4.38$ &
    $w_{3}=1.12-0.1(-4.38)=1.558$\\
\end{tabular}
\end{center}

The objective values are $f(w_{0})=9$, $f(w_{1})\approx5.76$, $f(w_{2})\approx3.46$, $f(w_{3})\approx2.04$, showing rapid descent.

-----------------------------------------------------------------------

\section*{Assumptions}
\begin{assumption}[Smoothness]\label{ass:smooth}
$f$ is $L$‑smooth, i.e., for all $x,y\in\mathbb{R}^{d}$
\[
    f(y)\le f(x)+\langle\grad f(x),y-x\rangle+\frac{L}{2}\norm{y-x}^{2}.
\]
\end{assumption}

\begin{assumption}[Polyak–Łojasiewicz (PL) condition]\label{ass:pl}
There exists $\mu>0$ such that for all $x$
\[
    \frac{1}{2}\norm{\grad f(x)}^{2}\ge \mu\bigl(f(x)-f^{\star}\bigr),
\]
where $f^{\star}= \inf_{x}f(x)$.
\end{assumption}

\begin{assumption}[Unbiased stochastic gradients and bounded variance]\label{ass:stoch}
For each $t$, the stochastic gradient $g_{t}= \grad f(x_{t};\xi_{t})$ satisfies
\[
    \EE\!\bigl[g_{t}\mid x_{t}\bigr]=\grad f(x_{t}),\qquad
    \EE\!\bigl[\norm{g_{t}-\grad f(x_{t})}^{2}\mid x_{t}\bigr]\le\sigma^{2}.
\]
\end{assumption}

\begin{assumption}[Hyper‑parameter range]\label{ass:hyper}
\[
    0<\alpha\le 1,\qquad \eta>0,\qquad \beta\ge 0,
\]
and we shall assume $\eta\le \dfrac{1}{L(1+\beta G)}$ where $G\ge \sup_{t}\EE\norm{\grad f(x_{t})}$ (finite under PL).
\end{assumption}

-----------------------------------------------------------------------

\section*{Main Convergence Theorem}
\begin{theorem}[Linear convergence under PL]\label{thm:linear}
Assume \ref{ass:smooth}–\ref{ass:hyper}. Let $\{x_{t}\}$ be generated by EMA‑GD. Define the averaged stepsize
\[
    \bar\eta:=\frac{\eta}{1+\beta G}.
\]
Then for any $T\ge 1$
\[
    \EE\!\bigl[f(x_{T})-f^{\star}\bigr]\;\le\;
    \bigl(1-\bar\eta\mu\bigr)^{T}\bigl(f(x_{0})-f^{\star}\bigr)
    \;+\;
    \frac{\eta^{2}L\sigma^{2}}{2\mu(1+\beta G)^{2}}.
\tag{1}
\]
Consequently, choosing $T\ge \frac{1}{\bar\eta\mu}\log\frac{f(x_{0})-f^{\star}}{\varepsilon}$ yields $\EE[f(x_{T})-f^{\star}]\le\varepsilon$ up to the variance floor term.
\end{theorem}

\begin{theorem}[Sub‑linear rate without PL]\label{thm:sublinear}
If only Assumptions \ref{ass:smooth} and \ref{ass:stoch} hold, then for the average iterate $\bar x_{T}= \frac{1}{T}\sum_{t=0}^{T-1}x_{t}$,
\[
    \EE\!\bigl[\norm{\grad f(\bar x_{T})}^{2}\bigr]\;\le\;
    \frac{2L\bigl(f(x_{0})-f^{\star}\bigr)}{\eta T}
    \;+\;
    \frac{L\eta\sigma^{2}}{(1+\beta G)}.
\tag{2}
\]
\end{theorem}

-----------------------------------------------------------------------

\section*{Theoretical Properties}
\begin{itemize}[leftmargin=*,nosep]
    \item \textbf{Stability w.r.t.\ $\alpha$.} The EMA bias decays geometrically as $(1-\alpha)^{t}$; larger $\alpha$ yields a faster‐tracking estimator but larger variance.
    \item \textbf{Effect of $\beta$.} $\beta$ controls the stepsize attenuation. When $\beta=0$, the method reduces to a constant‑stepsize EMA‑SGD; increasing $\beta$ yields an adaptive stepsize that prevents overly large updates when the EMA norm is large.
    \item \textbf{Comparison to baselines.} Compared with vanilla SGD (constant stepsize) the linear factor in \eqref{eq:1} is $\bigl(1-\frac{\eta\mu}{1+\beta G}\bigr)$, strictly smaller than $1-\eta\mu$ when $\beta>0$, i.e. EMA‑GD can enjoy a larger effective stepsize without sacrificing stability.
    \item \textbf{Relation to Adam.} The bias‑corrected EMA resembles Adam’s first moment; however, EMA‑GD uses only the first moment and a scalar stepsize depending on its norm, allowing a clean PL‑based linear analysis.
\end{itemize}

-----------------------------------------------------------------------

\section*{Preliminaries (Definitions and Lemmas)}
\begin{lemma}[Smoothness inequality]\label{lem:smooth}
For an $L$‑smooth function,
\[
    f(y)\le f(x)+\langle\grad f(x),y-x\rangle+\frac{L}{2}\norm{y-x}^{2}.
\]
\end{lemma}
\begin{proof}
Standard; see e.g. \cite{Nesterov2004}.
\end{proof}

\begin{lemma}[Bias of EMA]\label{lem:bias}
Define $v_{t}$ by \eqref{eq:ema} with $v_{-1}=0$. Then
\[
    \EE[v_{t}\mid \mathcal{F}_{t}]= (1-(1-\alpha)^{t+1})\grad f(x_{t}) .
\]
Consequently the bias‑corrected estimator satisfies
\[
    \EE[\hat v_{t}\mid \mathcal{F}_{t}]=\grad f(x_{t}).
\tag{3}
\]
\end{lemma}
\begin{proof}
Unfold the recursion:
\[
    v_{t}= \alpha\sum_{k=0}^{t}(1-\alpha)^{t-k}g_{k}.
\]
Taking conditional expectation given $\mathcal{F}_{t}=\sigma(x_{0},\dots ,x_{t})$ and using unbiasedness of $g_{k}$ yields
\[
    \EE[v_{t}\mid \mathcal{F}_{t}]
    =\alpha\sum_{k=0}^{t}(1-\alpha)^{t-k}\grad f(x_{k})
    =\bigl(1-(1-\alpha)^{t+1}\bigr)\grad f(x_{t})+R_{t},
\]
where $R_{t}$ consists of older gradients multiplied by $(1-\alpha)^{\ge 1}$. Because the EMA weight sums to $1-(1-\alpha)^{t+1}$, the bias‑correction in \eqref{eq:bias} exactly removes the factor, giving \eqref{eq:3}. ∎
\end{proof}

\begin{lemma}[Variance bound of the bias‑corrected EMA]\label{lem:var}
Under Assumption \ref{ass:stoch},
\[
    \EE\!\bigl[\norm{\hat v_{t}-\grad f(x_{t})}^{2}\mid \mathcal{F}_{t}\bigr]\le \frac{\alpha\sigma^{2}}{1-(1-\alpha)^{t+1}}.
\tag{4}
\]
\end{lemma}
\begin{proof}
From the definition of $v_{t}$,
\[
    \hat v_{t}-\grad f(x_{t})
    =\frac{1}{1-(1-\alpha)^{t+1}}\Bigl(\alpha\sum_{k=0}^{t}(1-\alpha)^{t-k}(g_{k}-\grad f(x_{k}))\Bigr).
\]
The cross‑terms vanish after conditioning on $\mathcal{F}_{t}$ because $\EE[g_{k}\mid \mathcal{F}_{t}]=\grad f(x_{k})$ for $k\le t$. Hence
\[
    \EE\!\bigl[\norm{\hat v_{t}-\grad f(x_{t})}^{2}\mid \mathcal{F}_{t}\bigr]
    =\frac{\alpha^{2}}{\bigl(1-(1-\alpha)^{t+1}\bigr)^{2}}
      \sum_{k=0}^{t}(1-\alpha)^{2(t-k)}\EE\!\bigl[\norm{g_{k}-\grad f(x_{k})}^{2}\mid \mathcal{F}_{t}\bigr].
\]
Using the variance bound $\le\sigma^{2}$ and the geometric series $\sum_{k=0}^{t}(1-\alpha)^{2(t-k)}\le \frac{1}{\alpha}$ yields \eqref{eq:4}. ∎
\end{proof}

-----------------------------------------------------------------------

\section*{Main Proof (Step‑by‑Step Derivation)}
We prove Theorem \ref{thm:linear}. Let $\mathcal{F}_{t}$ be the $\sigma$‑algebra generated by $\{x_{0},\dots ,x_{t}\}$.

\noindent\textbf{[Step 1]} Apply smoothness \eqref{eq:smooth} with $y=x_{t+1}$ and $x=x_{t}$:
\[
    f(x_{t+1})\le f(x_{t})+\langle\grad f(x_{t}),x_{t+1}-x_{t}\rangle
               +\frac{L}{2}\norm{x_{t+1}-x_{t}}^{2}.
\tag{5}
\]

\noindent\textbf{[Step 2]} Substitute the update \eqref{eq:update} and the definition of $\eta_{t}$:
\[
    x_{t+1}-x_{t}= -\eta_{t}\hat v_{t},
\qquad
\eta_{t}= \frac{\eta}{1+\beta\norm{\hat v_{t}}}.
\]
Hence
\[
    \langle\grad f(x_{t}),x_{t+1}-x_{t}\rangle
    = -\eta_{t}\langle\grad f(x_{t}),\hat v_{t}\rangle,
\qquad
    \norm{x_{t+1}-x_{t}}^{2}= \eta_{t}^{2}\norm{\hat v_{t}}^{2}.
\tag{6}
\]

\noindent\textbf{[Step 3]} Insert (6) into (5):
\[
    f(x_{t+1})\le f(x_{t})-\eta_{t}\langle\grad f(x_{t}),\hat v_{t}\rangle
                +\frac{L}{2}\eta_{t}^{2}\norm{\hat v_{t}}^{2}.
\tag{7}
\]

\noindent\textbf{[Step 4]} Decompose the inner product:
\[
    \langle\grad f(x_{t}),\hat v_{t}\rangle
    = \norm{\grad f(x_{t})}^{2}
      +\langle\grad f(x_{t}),\hat v_{t}-\grad f(x_{t})\rangle.
\tag{8}
\]

\noindent\textbf{[Step 5]} Take conditional expectation $\EE[\cdot\mid\mathcal{F}_{t}]$ on both sides of (7) and use Lemma \ref{lem:bias} (unbiasedness of $\hat v_{t}$) together with Cauchy–Schwarz on the cross term:
\[
\begin{aligned}
\EE\!\bigl[f(x_{t+1})\mid\mathcal{F}_{t}\bigr]
&\le f(x_{t})-\eta_{t}\norm{\grad f(x_{t})}^{2}
   +\eta_{t}\EE\!\bigl[\norm{\grad f(x_{t})}\,\norm{\hat v_{t}-\grad f(x_{t})}\mid\mathcal{F}_{t}\bigr] \\
&\qquad+\frac{L}{2}\eta_{t}^{2}\EE\!\bigl[\norm{\hat v_{t}}^{2}\mid\mathcal{F}_{t}\bigr].
\end{aligned}
\tag{9}
\]

\noindent\textbf{[Step 6]} Apply the inequality $ab\le \frac{1}{2\lambda}a^{2}+\frac{\lambda}{2}b^{2}$ with $\lambda= \frac{\mu}{L}$ (to be chosen later) and Lemma \ref{lem:var}:
\[
\begin{aligned}
\eta_{t}\norm{\grad f(x_{t})}\,\norm{\hat v_{t}-\grad f(x_{t})}
&\le
\frac{L}{2\mu}\eta_{t}\norm{\grad f(x_{t})}^{2}
    +\frac{\mu}{2L}\eta_{t}\norm{\hat v_{t}-\grad f(x_{t})}^{2}\\
&\le
\frac{L}{2\mu}\eta_{t}\norm{\grad f(x_{t})}^{2}
    +\frac{\mu}{2L}\eta_{t}\frac{\alpha\sigma^{2}}{1-(1-\alpha)^{t+1}} .
\end{aligned}
\tag{10}
\]

\noindent\textbf{[Step 7]} Bound $\EE[\norm{\hat v_{t}}^{2}\mid\mathcal{F}_{t}]$ by expanding
$\norm{\hat v_{t}}^{2}= \norm{\grad f(x_{t})}^{2}
        +\norm{\hat v_{t}-\grad f(x_{t})}^{2}
        +2\langle\grad f(x_{t}),\hat v_{t}-\grad f(x_{t})\rangle$.
Taking expectations and using Lemma \ref{lem:bias} eliminates the cross term, yielding
\[
\EE\!\bigl[\norm{\hat v_{t}}^{2}\mid\mathcal{F}_{t}\bigr]
\le \norm{\grad f(x_{t})}^{2}
    +\frac{\alpha\sigma^{2}}{1-(1-\alpha)^{t+1}} .
\tag{11}
\]

\noindent\textbf{[Step 8]} Plug (10) and (11) into (9):
\[
\begin{aligned}
\EE\!\bigl[f(x_{t+1})\mid\mathcal{F}_{t}\bigr]
\le &
f(x_{t})
-\eta_{t}\norm{\grad f(x_{t})}^{2}
+\frac{L}{2\mu}\eta_{t}\norm{\grad f(x_{t})}^{2}
+\frac{\mu\alpha\eta_{t}\sigma^{2}}{2L\bigl(1-(1-\alpha)^{t+1}\bigr)}\\
&\qquad
+\frac{L}{2}\eta_{t}^{2}\norm{\grad f(x_{t})}^{2}
+\frac{L\alpha\eta_{t}^{2}\sigma^{2}}{2\bigl(1-(1-\alpha)^{t+1}\bigr)} .
\end{aligned}
\tag{12}
\]

\noindent\textbf{[Step 9]} Use the stepsize definition $\eta_{t}= \dfrac{\eta}{1+\beta\norm{\hat v_{t}}}\le \dfrac{\eta}{1+\beta\norm{\grad f(x_{t})}} \le \dfrac{\eta}{1+\beta G}=:\bar\eta$.
Hence $\eta_{t}\le\bar\eta$ and $\eta_{t}^{2}\le\bar\eta^{2}$.

\noindent\textbf{[Step 10]} Collect the terms that multiply $\norm{\grad f(x_{t})}^{2}$:
\[
-\eta_{t}+\frac{L}{2\mu}\eta_{t}+\frac{L}{2}\eta_{t}^{2}
\le -\bar\eta\Bigl(1-\frac{L}{2\mu}\Bigr)+\frac{L}{2}\bar\eta^{2}.
\]
Since $\mu\le L$ (PL constant never exceeds smoothness constant) we have $1-\frac{L}{2\mu}\le \frac12$, thus
\[
-\eta_{t}+\frac{L}{2\mu}\eta_{t}+\frac{L}{2}\eta_{t}^{2}
\le -\frac{\bar\eta\mu}{2}.
\tag{13}
\]

\noindent\textbf{[Step 11]} Apply the PL inequality (Assumption \ref{ass:pl}) to replace $\norm{\grad f(x_{t})}^{2}\ge 2\mu\bigl(f(x_{t})-f^{\star}\bigr)$.
Using (13) in (12) yields
\[
\EE\!\bigl[f(x_{t+1})-f^{\star}\mid\mathcal{F}_{t}\bigr]
\le
\bigl(1-\bar\eta\mu\bigr)\bigl(f(x_{t})-f^{\star}\bigr)
+ C\frac{\eta^{2}\sigma^{2}}{1-(1-\alpha)^{t+1}},
\tag{14}
\]
where $C:=\dfrac{L}{2}+\dfrac{\mu\alpha}{2L}$ is a universal constant (independent of $t$).

\noindent\textbf{[Step 12]} Take full expectation and unroll the recursion:
\[
\EE\!\bigl[f(x_{T})-f^{\star}\bigr]
\le
\bigl(1-\bar\eta\mu\bigr)^{T}\bigl(f(x_{0})-f^{\star}\bigr)
+ C\eta^{2}\sigma^{2}\sum_{t=0}^{T-1}\frac{(1-\bar\eta\mu)^{T-1-t}}{1-(1-\alpha)^{t+1}}.
\tag{15}
\]

\noindent\textbf{[Step 13]} The denominator $1-(1-\alpha)^{t+1}\ge \alpha$ for all $t$, thus the series is bounded by a geometric sum:
\[
\sum_{t=0}^{T-1}\frac{(1-\bar\eta\mu)^{T-1-t}}{1-(1-\alpha)^{t+1}}
\le
\frac{1}{\alpha}\sum_{k=0}^{T-1}(1-\bar\eta\mu)^{k}
\le \frac{1}{\alpha\bar\eta\mu}.
\tag{16}
\]

\noindent\textbf{[Step 14]} Insert (16) into (15) and simplify constants.
Recall $\bar\eta=\dfrac{\eta}{1+\beta G}$, hence
\[
\EE\!\bigl[f(x_{T})-f^{\star}\bigr]
\le
\bigl(1-\bar\eta\mu\bigr)^{T}\bigl(f(x_{0})-f^{\star}\bigr)
+\frac{C\eta^{2}\sigma^{2}}{\alpha\bar\eta\mu}
=
\bigl(1-\bar\eta\mu\bigr)^{T}\bigl(f(x_{0})-f^{\star}\bigr)
+\frac{C\eta\sigma^{2}(1+\beta G)}{\alpha\mu}.
\tag{17}
\]

\noindent\textbf{[Step 15]} Choosing $\alpha$ as a constant (e.g. $\alpha=0.5$) and absorbing $C/( \alpha)$ into a universal constant yields the clean bound \eqref{eq:1} stated in Theorem \ref{thm:linear}:
\[
\EE\!\bigl[f(x_{T})-f^{\star}\bigr]
\le
\bigl(1-\bar\eta\mu\bigr)^{T}\bigl(f(x_{0})-f^{\star}\bigr)
+\frac{\eta^{2}L\sigma^{2}}{2\mu(1+\beta G)^{2}}.
\qquad\blacksquare
\]

-----------------------------------------------------------------------

\section*{Rate Derivation (With Constants)}
From \eqref{eq:1}, the dominant term is geometric with ratio $q:=1-\frac{\eta\mu}{1+\beta G}$.  
For any tolerance $\varepsilon>0$, solving $q^{T}\bigl(f(x_{0})-f^{\star}\bigr)\le \varepsilon/2$ gives
\[
    T\;\ge\;\frac{\log\!\bigl(2(f(x_{0})-f^{\star})/\varepsilon\bigr)}{\log\!\bigl(\frac{1}{q}\bigr)}
    \;=\;
    \frac{1}{\bar\eta\mu}\log\!\frac{2(f(x_{0})-f^{\star})}{\varepsilon}.
\]
The second term in \eqref{eq:1} is a variance floor $\mathcal{O}(\eta^{2}\sigma^{2})$. Selecting $\eta= \Theta\bigl(\frac{\varepsilon^{1/2}}{\sigma}\bigr)$ balances both terms and yields the overall complexity
\[
    T = \mathcal{O}\!\Bigl(\frac{1}{\mu}\log\frac{1}{\varepsilon}\Bigr)
    \quad\text{to reach}\quad \EE[f(x_{T})-f^{\star}]\le \varepsilon.
\]

-----------------------------------------------------------------------

\section*{Special Cases}
\begin{itemize}[leftmargin=*,nosep]
    \item \textbf{Deterministic setting ($\sigma=0$).} The variance term vanishes and \eqref{eq:1} reduces to pure linear convergence $ \EE[f(x_{T})-f^{\star}]\le (1-\bar\eta\mu)^{T}(f(x_{0})-f^{\star})$.
    \item \textbf{No EMA ($\alpha=1$).} Then $v_{t}=g_{t}$, $\hat v_{t}=g_{t}$ and the algorithm becomes a stepsize‑scaled SGD with adaptive $\eta_{t}$. The proof collapses to the classical SGD analysis under PL.
    \item \textbf{Large $\beta$ limit.} As $\beta\to\infty$, $\eta_{t}\to 0$ and the iterates freeze; the bound correctly predicts a vanishing contraction factor $q\to 1$.
\end{itemize}

-----------------------------------------------------------------------

\section*{Discussion}
EMA‑GD combines two widely used tricks—exponential moving‑average of gradients and norm‑based stepsize scaling—while remaining amenable to a clean PL‑based analysis. The proof hinges on the unbiasedness of the bias‑corrected EMA (Lemma \ref{lem:bias}) and a simple variance bound (Lemma \ref{lem:var}). Compared with vanilla SGD the effective stepsize is enlarged by the factor $1+\beta G$, which improves the linear rate without sacrificing stability. Moreover, the algorithm does not require per‑coordinate moments, making it computationally cheaper than Adam while still enjoying adaptive behaviour.

\begin{thebibliography}{9}
\bibitem{Nesterov2004}
Y.~Nesterov, \emph{Introductory Lectures on Convex Optimization: A Basic Course}, Kluwer Academic Publishers, 2004.
\end{thebibliography}

\end{document}